In the majority of this chapter, we will consider sequences $a_n$ defined
\emph{by recurrence} or \emph{recursively} by the conditions:
\index{sequence!recursive}
\index{sequence!defined by recurrence}
\index{recursion}
\begin{equation}\label{eq1}
\begin{cases}
  a_0 = \alpha,\\
  a_{n+1} = f(a_n)
\end{cases}
\end{equation}

Given the initial term $\alpha$ and the recurrence law $f$, there is a unique
sequence that satisfies \eqref{eq1}, and its terms are:
\begin{align*}
a_0 & =\alpha,\\
a_1 &= f(a_0)=f(\alpha),\\
a_2 &= f(a_1)=f(f(\alpha)),\\
a_3 &= f(a_2)=f(f(f(\alpha))),\\
&\ \vdots\\
a_n &= f(a_{n-1}) = f^n(\alpha),\\
&\ \vdots
\end{align*}

The value of $a_n$ might represent the state of a system that evolves from an
initial state $a_0=\alpha$ through the function $f$, which represents the change
of state. The natural number $n$ could thus represent a time step. In this
sense, \eqref{eq1} is also called a \emph{discrete dynamical system}%
\mymargin{discrete dynamical system}\index{system!dynamical!discrete}.

The equation $a_{n+1} = f(a_n)$ is called an \emph{autonomous first-order
recursive equation}. There are other types of equations that we will consider
only marginally in the following chapters. For example, when we defined the
factorial: $a_n = n!$, we provided the conditions:
\[
\begin{cases}
  a_0 = 1\\
  a_{n+1} = (n+1) \cdot a_n
\end{cases}
\]
but the equation $a_{n+1} = (n+1) \cdot a_n$ is of the form $a_{n+1} = f(n,
a_n)$ and is said to be \emph{non-autonomous} because the recurrence function
$f$ explicitly depends on $n$ as well as on the previous term $a_n$.

Equations of order higher than the first could also be considered. For example,
the sequence $F_n$ of \emph{Fibonacci} (Leonardo Pisano, 1170--1242), whose
first terms are listed in Table~\ref{tab:Fibonacci}
\begin{table}
  \begin{center}
  0, 1, 1, 2, 3, 5, 8, 13, 21, 34, 55, 89, 144,
  233, 377, 610, 987\dots
\end{center}
\caption{The first terms of the Fibonacci sequence. Each term is the sum of the two preceding ones.}
\label{tab:Fibonacci}
\end{table}
\mymargin{Fibonacci}%
\index{Fibonacci}%
\index{Fibonacci!sequence of}%
\index{sequence!Fibonacci}%
satisfies the recursive equation:
\begin{equation}\label{eq:Fibonacci}
\begin{cases}
  F_0 = 0 \\
  F_1 = 1 \\
  F_{n+2} = F_{n+1} + F_n
\end{cases}
\end{equation}
which is a second-order relation since each term can be defined using the values
of the \emph{two} preceding terms. If the equation is linear, as in this case,
explicit formulas can be found to express the $n$-th term of the sequence. We
will do this in section~\ref{sec:ricorrenza_lineare}.

Systems of recursive equations could also be considered. For example, if $f$
were a complex function $f\colon \CC \to \CC$, $f(x+iy) = f_1(x,y) + i f_2(x,y)$
with $f_1,f_2 \colon \RR\times\RR\to\RR$, one could write $a_n = x_n + i y_n$
with $x_n, y_n \in \RR$, and the recursive equation $a_{n+1} = f(a_n)$ would
become a system of two equations:
\[
  \begin{cases}
    x_{n+1} = f_1(x_n, y_n)\\
    y_{n+1} = f_2(x_n, y_n).
  \end{cases}
\]
The study of systems goes beyond the scope of this chapter, but we will briefly
mention an example in section~\ref{sec:mandelbrot}.